% =========================================================================
% COVER PAGE / OUTER FRONT PAGE
% =========================================================================
\thispagestyle{empty}
\begin{center}
    {\large \textbf{VISVESVARAYA TECHNOLOGICAL UNIVERSITY}}\\
    {\normalsize \textbf{Jnana Sangama, Machhe, Belagavi - 590018, Karnataka, India}}\\[0.2cm]

    \includegraphics[width=2.5cm]{vtu_logo.png}\\[0.2cm]

    {\large \textbf{A PROJECT REPORT (2022 SCHEME)}}\\
    {\normalsize \textbf{ON}}\\[0.12cm]

    {\Large \textbf{``AI-BASED SNOOKER BALL IDENTIFICATION AND LASER POSITIONING SYSTEM''}}\\[0.2cm]

    {\small \textit{Submitted in partial fulfillment of the academic requirements for the award of the degree of}}\\[0.08cm]
    {\large \textbf{BACHELOR OF ENGINEERING}}\\[0.04cm]
    {\normalsize \textbf{IN}}\\[0.04cm]
    {\large \textbf{ARTIFICIAL INTELLIGENCE AND MACHINE LEARNING}}\\[0.2cm]

    \textbf{Submitted by}\\[0.08cm]
    \begin{tabular}{ll}
        \textbf{AAMIR AWAIZ PASHA} & \textbf{(4PM23AI001)} \\
        \textbf{AYAN KHAN}         & \textbf{(4PM23AI005)} \\
        \textbf{IBAD SHARIEF}       & \textbf{(4PM23AI016)} \\
        \textbf{MOHAMMED OWAIES}   & \textbf{(4PM23AI026)} \\
    \end{tabular}

    \vspace{0.15cm}
    {\small \textit{Under the Guidance of}}\\[0.06cm]
    {\textbf{Mrs. Swetha B S}}\\[0.02cm]
    {\small Assistant Professor, Dept. of AIML, PESITM, Shivamogga}\\[0.2cm]

    \includegraphics[width=4.6cm]{college_logo.jpg}\\[0.15cm]

    {\large \textbf{DEPARTMENT OF ARTIFICIAL INTELLIGENCE AND MACHINE LEARNING}}\\[0.04cm]
    {\textbf{P E S INSTITUTE OF TECHNOLOGY AND MANAGEMENT}}\\[0.03cm]
    {\small (Affiliated to Visvesvaraya Technological University, Belagavi, Approved by AICTE, New Delhi)}\\
    {\small NH-206, Sagar Road, Shivamogga - 577204, Karnataka, India}\\
    {\textbf{2026 -- 2027}}
\end{center}

% =========================================================================
% INNER TITLE PAGE
% =========================================================================
\newpage
\thispagestyle{empty}
\begin{center}
    {\large \textbf{VISVESVARAYA TECHNOLOGICAL UNIVERSITY}}\\
    {\normalsize \textbf{Jnana Sangama, Machhe, Belagavi - 590018, Karnataka, India}}\\[0.2cm]

    \includegraphics[width=2.5cm]{vtu_logo.png}\\[0.2cm]

    {\large \textbf{A PROJECT REPORT (2022 SCHEME)}}\\
    {\normalsize \textbf{ON}}\\[0.12cm]

    {\Large \textbf{``AI-BASED SNOOKER BALL IDENTIFICATION AND LASER POSITIONING SYSTEM''}}\\[0.2cm]

    {\small \textit{Submitted in partial fulfillment of the academic requirements for the award of the degree of}}\\[0.08cm]
    {\large \textbf{BACHELOR OF ENGINEERING}}\\[0.04cm]
    {\normalsize \textbf{IN}}\\[0.04cm]
    {\large \textbf{ARTIFICIAL INTELLIGENCE AND MACHINE LEARNING}}\\[0.2cm]

    \textbf{Submitted by}\\[0.08cm]
    \begin{tabular}{ll}
        \textbf{AAMIR AWAIZ PASHA} & \textbf{(4PM23AI001)} \\
        \textbf{AYAN KHAN}         & \textbf{(4PM23AI005)} \\
        \textbf{IBAD SHARIEF}       & \textbf{(4PM23AI016)} \\
        \textbf{MOHAMMED OWAIES}   & \textbf{(4PM23AI026)} \\
    \end{tabular}

    \vspace{0.15cm}
    {\small \textit{Under the Guidance of}}\\[0.06cm]
    {\textbf{Mrs. Swetha B S}}\\[0.02cm]
    {\small Assistant Professor, Dept. of AIML, PESITM, Shivamogga}\\[0.2cm]

    \includegraphics[width=4.6cm]{college_logo.jpg}\\[0.15cm]

    {\large \textbf{DEPARTMENT OF ARTIFICIAL INTELLIGENCE AND MACHINE LEARNING}}\\[0.04cm]
    {\textbf{P E S INSTITUTE OF TECHNOLOGY AND MANAGEMENT}}\\[0.03cm]
    {\small (Affiliated to Visvesvaraya Technological University, Belagavi, Approved by AICTE, New Delhi)}\\
    {\small NH-206, Sagar Road, Shivamogga - 577204, Karnataka, India}\\
    {\textbf{2026 -- 2027}}
\end{center}

% =========================================================================
% VISION, MISSION, PEOs, PSOs, AND POs (PESITM DEPT OF AIML)
% =========================================================================
\newpage
\thispagestyle{frontmatterstyle}
\addcontentsline{toc}{chapter}{VISION, MISSION, PEOs, PSOs AND POs}

\noindent{\large \textbf{\underline{Institute Vision}}}\\
To be the most preferred institution for engineering \& management education, research \& entrepreneurship by creating professionally superior \& ethically strong global manpower.

\vspace{0.3cm}
\noindent{\large \textbf{\underline{Institute Mission}}}\\
\begin{enumerate}
    \item Provide high-quality engineering and management education with technical excellence for global requirements.
    \item Promote research and entrepreneurship with collaboration, modern facilities, and strong mentorship.
    \item Build ethically grounded graduates who succeed globally, lead with integrity and commit to lifelong learning and social responsibility.
\end{enumerate}

\vspace{0.3cm}
\noindent{\large \textbf{\underline{Department Vision}}}\\
To be the frontiers of AIML and to nurture ethically grounded professionals who can contribute innovatively to society while cultivating global competency through research, entrepreneurship and innovation.

\vspace{0.3cm}
\noindent{\large \textbf{\underline{Department Mission}}}\\
\begin{itemize}
    \item \textbf{Mission 1:} Provide top-notch AIML education that includes ethics, preparing professionals to use technology for good.
    \item \textbf{Mission 2:} Support innovative and entrepreneurial thinking in AIML, turning ideas into impactful solutions.
    \item \textbf{Mission 3:} Promote teamwork across disciplines and cultures to solve global challenges with AIML.
\end{itemize}

\newpage
\thispagestyle{frontmatterstyle}

\noindent{\large \textbf{\underline{Program Educational Objectives (PEOs)}}}\\
\begin{itemize}
    \item \textbf{PEO 1:} To conceptualize, analyze, design and develop IT Solutions of varying complexities by leveraging advances in computer technology.
    \item \textbf{PEO 2:} To apply standard practices and strategies in software project development and management using industry-wide benchmarked framework to deliver a sustainable quality product.
    \item \textbf{PEO 3:} To work as a team player in a cross-cultural environment with a passion for entrepreneurship and a zest for higher studies.
\end{itemize}

\vspace{0.4cm}
\noindent{\large \textbf{\underline{Program Specific Outcomes (PSOs)}}}\\
\begin{itemize}
    \item \textbf{PSO 1:} Build skills to develop software applications in specialized areas of Computer Science and Engineering such as Artificial Intelligence, Machine Learning, Data Science, Web Development, Gaming, Augmented Reality / Virtual Reality (AR/VR).
    \item \textbf{PSO 2:} Focus on exploring supervised, unsupervised and reinforcement learning and apply them to a range of AI problems.
    \item \textbf{PSO 3:} Develop an understanding of the evolution of AI and future developments in the AI revolution.
\end{itemize}

\newpage
\thispagestyle{frontmatterstyle}

\noindent{\large \textbf{\underline{Program Outcomes (POs)}}}\\
\begin{spacing}{1.12}
{\small
\begin{itemize}
    \item \textbf{PO1: Engineering knowledge:} Apply the knowledge of mathematics, science, engineering fundamentals, and an engineering specialization for the solution of complex engineering problems.
    \item \textbf{PO2: Problem analysis:} Identify, formulate, research literature, and analyze complex engineering problems reaching substantiated conclusions using first principles of mathematics, natural sciences, and engineering sciences.
    \item \textbf{PO3: Design/development of solutions:} Design solutions for complex engineering problems and design system components or processes that meet the specified needs with appropriate consideration for public health and safety and cultural, societal, and environmental considerations.
    \item \textbf{PO4: Conduct investigations of complex problems:} Use research-based knowledge and research methods including design of experiments, analysis and interpretation of data, and synthesis of the information to provide valid conclusions.
    \item \textbf{PO5: Modern tool usage:} Create, select, and apply appropriate techniques, resources, and modern engineering and IT tools, including prediction and modeling to complex engineering activities, with an understanding of the limitations.
    \item \textbf{PO6: The engineer and society:} Apply reasoning informed by the contextual knowledge to assess societal, health, safety, legal, and cultural issues and the consequent responsibilities relevant to the professional engineering practice.
    \item \textbf{PO7: Environment and sustainability:} Understand the impact of the professional engineering solutions in societal and environmental contexts, and demonstrate the knowledge of, and need for sustainable development.
    \item \textbf{PO8: Ethics:} Apply ethical principles and commit to professional ethics and responsibilities and norms of the engineering practice.
    \item \textbf{PO9: Individual and team work:} Function effectively as an individual, and as a member or leader in diverse teams, and in multidisciplinary settings.
    \item \textbf{PO10: Communication:} Communicate effectively on complex engineering activities with the engineering community and with the society at large, such as, being able to comprehend and write effective reports and design documentation, make effective presentations, and give and receive clear instructions.
    \item \textbf{PO11: Project management and finance:} Demonstrate knowledge and understanding of the engineering and management principles and apply these to one's own work, as a member and leader in a team, to manage projects and in multidisciplinary environments.
    \item \textbf{PO12: Life-long learning:} Recognize the need for, and have the preparation and ability to engage in independent and life-long learning in the broadest context of technological change.
\end{itemize}
}
\end{spacing}

% =========================================================================
% CERTIFICATE PAGE
% =========================================================================
\newpage
\thispagestyle{empty}
\addcontentsline{toc}{chapter}{CERTIFICATE}
\begin{center}
    \includegraphics[width=3.6cm]{college_logo.jpg}\\[0.08cm]
    {\large \textbf{P E S INSTITUTE OF TECHNOLOGY AND MANAGEMENT}}\\[0.03cm]
    {\small NH-206, Sagar Road, Shivamogga - 577204, Karnataka, India}\\
    {\small (Affiliated to Visvesvaraya Technological University, Belagavi)}\\[0.08cm]
    {\large \textbf{DEPARTMENT OF ARTIFICIAL INTELLIGENCE AND MACHINE LEARNING}}\\[0.12cm]
    {\Large \textbf{\underline{CERTIFICATE}}}\\[0.15cm]
\end{center}

\noindent \begin{spacing}{1.15}
{\small This is to certify that the project work entitled \textbf{``AI-BASED SNOOKER BALL IDENTIFICATION AND LASER POSITIONING SYSTEM''} has been successfully carried out by \textbf{Aamir Awaiz Pasha (4PM23AI001)}, \textbf{Ayan Khan (4PM23AI005)}, \textbf{Ibad Sharief (4PM23AI016)}, and \textbf{Mohammed Owaies (4PM23AI026)}, bonafide students of \textbf{PES Institute of Technology and Management, Shivamogga}, in partial fulfillment for the award of the degree of \textbf{Bachelor of Engineering in Artificial Intelligence and Machine Learning} of \textbf{Visvesvaraya Technological University, Belagavi} during the academic year \textbf{2026--2027}.

It is certified that all corrections and suggestions indicated during internal evaluations have been fully incorporated in the project report deposited in the departmental library. The project report has been approved as it satisfies the academic requirements in respect of project work prescribed for the said degree.}
\end{spacing}

\vspace{0.25cm}
\noindent
\begin{tabular}{cc}
    \rule{5.8cm}{0.4pt} & \rule{5.8cm}{0.4pt} \\
    \textbf{Mrs. Swetha B S} & \textbf{Mrs. Poornima B P} \\
    Assistant Professor & Assistant Professor \\
    Project Guide & Project Coordinator \\
    Dept. of AIML, PESITM & Dept. of AIML, PESITM \\[0.45cm]
    \rule{5.8cm}{0.4pt} & \rule{5.8cm}{0.4pt} \\
    \textbf{Dr. Likewin Thomas} & \textbf{Dr. Swamy D. R.} \\
    Professor and HOD & Principal \\
    Dept. of AIML, PESITM & PESITM, Shivamogga \\
\end{tabular}

\vspace{0.2cm}
\begin{center}
    \textbf{\underline{EXTERNAL VIVA EVALUATION}}
\end{center}
\vspace{0.05cm}
\noindent
\begin{tabular}{|c|p{7.6cm}|p{4.8cm}|}
    \hline
    \textbf{Sl. No.} & \textbf{Name of External / Internal Examiner} & \textbf{Signature with Date} \\
    \hline
    \rule{0pt}{0.55cm} 1. & & \\[0.2cm]
    \hline
    \rule{0pt}{0.55cm} 2. & & \\[0.2cm]
    \hline
\end{tabular}

% =========================================================================
% STUDENT DECLARATION
% =========================================================================
\newpage
\thispagestyle{frontmatterstyle}
\addcontentsline{toc}{chapter}{DECLARATION}
\begin{center}
    {\Large \textbf{\underline{DECLARATION}}}\\[0.8cm]
\end{center}

\noindent We, the undersigned students of the Department of Artificial Intelligence and Machine Learning, \textbf{PES Institute of Technology and Management, Shivamogga}, hereby declare that the project work entitled \textbf{``AI-BASED SNOOKER BALL IDENTIFICATION AND LASER POSITIONING SYSTEM''} has been independently executed by us under the guidance of our Project Guide, \textbf{Mrs. Swetha B S}, Assistant Professor, Department of AI \& ML, and Project Coordinator, \textbf{Mrs. Poornima B P}, Assistant Professor, Department of AI \& ML, PESITM, Shivamogga, in partial fulfillment of the requirements for the award of the degree of \textbf{Bachelor of Engineering in Artificial Intelligence and Machine Learning} by \textbf{Visvesvaraya Technological University, Belagavi}, during the academic year \textbf{2026--2027}.

\noindent We further declare that this report has not previously formed the basis for the award of any Degree, Diploma, Associateship, Fellowship, or any other similar title to any university or institution. We have properly acknowledged and cited all literature, open-source code libraries, algorithms, and third-party resources utilized throughout this work in accordance with academic integrity guidelines.

\vspace{1.2cm}
\noindent
\textbf{Place: Shivamogga}\\
\textbf{Date: \today}

\vspace{0.8cm}
\noindent
\begin{tabular}{ll}
    \rule{5.8cm}{0.4pt} & \rule{5.8cm}{0.4pt} \\
    \textbf{AAMIR AWAIZ PASHA} & \textbf{AYAN KHAN} \\
    (USN: 4PM23AI001) & (USN: 4PM23AI005) \\[1.2cm]
    \rule{5.8cm}{0.4pt} & \rule{5.8cm}{0.4pt} \\
    \textbf{IBAD SHARIEF} & \textbf{MOHAMMED OWAIES} \\
    (USN: 4PM23AI016) & (USN: 4PM23AI026) \\
\end{tabular}

% =========================================================================
% ACKNOWLEDGEMENTS
% =========================================================================
\newpage
\thispagestyle{frontmatterstyle}
\addcontentsline{toc}{chapter}{ACKNOWLEDGEMENT}
\begin{center}
    {\Large \textbf{\underline{ACKNOWLEDGEMENT}}}\\[0.6cm]
\end{center}

\noindent The satisfaction and euphoria that accompany the successful completion of any academic endeavor would remain incomplete without the heartfelt expression of gratitude to the esteemed individuals whose guidance, encouragement, and support made this project possible.

\noindent First and foremost, we express our profound gratitude and reverence to the \textbf{Management of PES Trust}, Shivamogga, for providing outstanding academic infrastructure, high-performance computing laboratories, and an inspiring academic environment.

\noindent We convey our heartfelt thanks to our respected Principal, \textbf{Dr. Swamy D. R.}, for his visionary leadership, continuous encouragement, and for providing the necessary institutional facilities to execute our experimental mechatronics setup.

\noindent We express our deep sense of gratitude to our respected Head of the Department, \textbf{Dr. Likewin Thomas}, Professor and HOD, Department of Artificial Intelligence and Machine Learning, for his constant academic encouragement, constructive suggestions, and relentless motivation throughout our undergraduate curriculum.

\noindent We express our profound debt of gratitude to our respected \textbf{Project Guide}, \textbf{Mrs. Swetha B S}, Assistant Professor, Department of Artificial Intelligence and Machine Learning, for her invaluable supervision, insightful technical advice, patient reviews, and unfailing encouragement from the inception to the successful completion of this project.

\noindent We convey our sincere gratitude to our respected \textbf{Project Coordinator}, \textbf{Mrs. Poornima B P}, Assistant Professor, Department of Artificial Intelligence and Machine Learning, for her constant coordination, timely advice, and diligent tracking of our project milestones.

\noindent We also take this opportunity to express our sincere appreciation to the \textbf{Project Review Committee (PRC)} members and all faculty and non-teaching staff members of the Department of Artificial Intelligence and Machine Learning for their kind cooperation, technical guidance, and valuable constructive critiques during periodic review evaluations.

\noindent Finally, our heartfelt gratitude goes out to our beloved \textbf{Parents and Family Members} for their unconditional love, moral support, and endless sacrifices. We also thank our fellow batchmates and friends who assisted directly or indirectly during physical gantry fabrication, camera calibration trials, and hardware testing.

\vspace{0.8cm}
\noindent
\begin{flushright}
    \textbf{Aamir Awaiz Pasha (4PM23AI001)}\\
    \textbf{Ayan Khan (4PM23AI005)}\\
    \textbf{Ibad Sharief (4PM23AI016)}\\
    \textbf{Mohammed Owaies (4PM23AI026)}\\
\end{flushright}

% =========================================================================
% ABSTRACT
% =========================================================================
\newpage
\thispagestyle{frontmatterstyle}
\addcontentsline{toc}{chapter}{ABSTRACT}
\begin{center}
    {\Large \textbf{\underline{ABSTRACT}}}\\[0.4cm]
\end{center}

\noindent In competitive snooker governed by the World Professional Billiards and Snooker Association (WPBSA), strict foul-and-a-miss regulations (Rule 14) frequently require disturbed balls to be restored to their exact pre-shot positions at the request of the non-offending player. In contemporary professional tournaments, this restoration relies almost entirely on visual estimation by human referees, occasionally assisted by broadcast split-screen television monitors, leading to notable positioning inaccuracies, player contention, and protracted match stoppages. In club matches and qualifying rounds, no broadcast infrastructure exists, leaving referees to replace balls purely from subjective human memory.

\noindent This dissertation presents the design, mathematical formulation, hardware fabrication, and software realization of an end-to-end, automated \textbf{AI-Based Snooker Ball Identification and Laser Positioning System}. The setup features an institutional $1168.4\,\text{mm} \times 2341.88\,\text{mm}$ table spanned by an overhead structural PVC gantry ($1315.15\,\text{mm}$ clearance) equipped with an overhead vision sensor and an IoT-enabled dual-axis micro-servo pan-tilt laser module controlled by an ESP32 microcontroller operating in dedicated Wi-Fi station mode (`10.21.46.190`). The perception pipeline integrates a YOLOv8s deep convolutional neural network fine-tuned on a custom snooker ball dataset comprising White, Pink, Blue, Green, Yellow, and Orange object balls. To overcome projective distortion and lens perspective warping, a planar homography matrix $\mathbf{H} \in \mathbb{R}^{3 \times 3}$ is computed via SVD across calibrated table fiducials, combined with a closed-form spherical elevation parallax correction for $52.5\,\text{mm}$ balls.

\noindent To eliminate micro-servo mechanical backlash and gravity droop, an empirical 2D Bivariate Cubic Spline interpolation model was calibrated across a 9-point spatial grid. The software system incorporates an autonomous referee engine that captures pre-shot baseline frames, tracks real-time motion via temporal frame differencing, automatically confirms post-shot settlement, isolates displaced balls via Hungarian bipartite matching, and drives the steerable laser diode to pinpoint exact restoration coordinates on the table baize. A full-stack web interface featuring a Three.js 3D interactive digital twin, live servo nudge fine-tuning, and offline-resilient frame caching was deployed. Experimental evaluations demonstrate a ball detection $\text{mAP}_{0.5}$ of $94.6\%$, an average homography metric localization error of $1.82\,\text{mm}$, and a laser targeting precision of $\pm 2.18\,\text{mm}$, with end-to-end analysis and actuation completing in $307\,\text{ms}$.

\vspace{0.4cm}
\noindent \textbf{Keywords:} Snooker, Computer Vision, Artificial Intelligence, YOLOv8, Planar Homography, Elevation Parallax, Bivariate Spline Interpolation, Dual-Axis Laser Steering, ESP32 Microcontroller, Autonomous Referee, Foul and a Miss, Ball Restoration.

% =========================================================================
% TABLE OF CONTENTS & LISTS
% =========================================================================
\newpage
\tableofcontents

\newpage
\addcontentsline{toc}{chapter}{LIST OF TABLES}
\listoftables

\newpage
\addcontentsline{toc}{chapter}{LIST OF FIGURES}
\listoffigures

% =========================================================================
% LIST OF ABBREVIATIONS
% =========================================================================
\newpage
\thispagestyle{frontmatterstyle}
\addcontentsline{toc}{chapter}{LIST OF ABBREVIATIONS}
\begin{center}
    {\Large \textbf{\underline{LIST OF ABBREVIATIONS}}}\\[0.6cm]
\end{center}

\begin{table}[H]
\centering
\begin{tabular}{ll}
    \textbf{AI}      & Artificial Intelligence \\
    \textbf{API}     & Application Programming Interface \\
    \textbf{BOM}     & Bill of Materials \\
    \textbf{CHT}     & Circular Hough Transform \\
    \textbf{CNN}     & Convolutional Neural Network \\
    \textbf{CPU}     & Central Processing Unit \\
    \textbf{ESP32}   & Espressif Systems 32-bit Microcontroller \\
    \textbf{FPS}     & Frames Per Second \\
    \textbf{GPIO}    & General Purpose Input/Output \\
    \textbf{GPU}     & Graphics Processing Unit \\
    \textbf{HOD}     & Head of the Department \\
    \textbf{HSV}     & Hue, Saturation, Value Color Space \\
    \textbf{HTTP}    & Hypertext Transfer Protocol \\
    \textbf{IoT}     & Internet of Things \\
    \textbf{JSON}    & JavaScript Object Notation \\
    \textbf{mAP}     & Mean Average Precision \\
    \textbf{NMS}     & Non-Maximum Suppression \\
    \textbf{OpenCV}  & Open Source Computer Vision Library \\
    \textbf{PVC}     & Polyvinyl Chloride \\
    \textbf{PWM}     & Pulse Width Modulation \\
    \textbf{REST}    & Representational State Transfer \\
    \textbf{RGB}     & Red, Green, Blue Color Space \\
    \textbf{RoI}     & Region of Interest \\
    \textbf{SVD}     & Singular Value Decomposition \\
    \textbf{Three.js}& 3D WebGL JavaScript Graphics Library \\
    \textbf{UI}      & User Interface \\
    \textbf{USN}     & University Seat Number \\
    \textbf{VTU}     & Visvesvaraya Technological University \\
    \textbf{WebGL}   & Web Graphics Library \\
    \textbf{WPBSA}   & World Professional Billiards and Snooker Association \\
    \textbf{YOLO}    & You Only Look Once \\
\end{tabular}
\end{table}
